% ============================================================
% BAGIAN 1 - ALGORITMA PATHFINDING
% ============================================================

Permasalahan \textit{Ice Sliding Puzzle Solver} pada tugas ini diselesaikan sebagai persoalan pencarian jalur pada ruang state. 
Berbeda dengan pathfinding biasa yang umumnya bergerak satu petak setiap langkah, pada puzzle ini aktor bergerak secara \textit{sliding}, yaitu terus berjalan pada satu arah sampai tepat sebelum menabrak rintangan. 
Oleh karena itu, satu aksi gerakan dapat melewati beberapa tile sekaligus dan cost gerakannya dihitung dari seluruh tile yang dilalui.

\begin{tcolorbox}[
	colback=blue!5,
	colframe=blue!60!black,
	title=\textbf{Inti Permasalahan Pathfinding},
	arc=2mm,
	boxrule=0.8pt
]
Algoritma harus mencari urutan gerakan dari posisi awal \texttt{Z} menuju titik tujuan \texttt{O}, dengan syarat aktor harus melewati angka secara berurutan dari \texttt{0} sampai angka terbesar yang ada pada papan. 
Solusi hanya dianggap valid apabila aktor berhenti tepat di \texttt{O} dan seluruh angka yang diperlukan sudah dilewati sesuai urutan.
\end{tcolorbox}

\section{Representasi State}

Setiap konfigurasi pencarian direpresentasikan sebagai sebuah state. 
State tidak cukup hanya menyimpan posisi aktor, karena permainan juga memiliki constraint angka berurutan. 
Oleh karena itu, state juga menyimpan informasi angka yang sudah dikumpulkan, angka terakhir yang valid, path gerakan, serta cost yang sudah dikeluarkan.

\begin{table}[H]
	\centering
	\caption{Representasi state pada algoritma pathfinding}
	\begin{tabular}{|p{4cm}|p{10cm}|}
		\hline
		\rowcolor{blue!15}
		\textbf{Komponen State} & \textbf{Keterangan} \\
		\hline
		Posisi aktor & Menyimpan koordinat baris dan kolom aktor pada papan. \\
		\hline
		\textit{Collected mask} & Menyimpan himpunan angka yang sudah dikumpulkan dalam bentuk bitmask. \\
		\hline
		Angka terakhir & Menyimpan angka terakhir yang berhasil dilewati secara valid. Nilai awalnya adalah $-1$. \\
		\hline
		Path gerakan & Menyimpan urutan gerakan yang sudah dilakukan, misalnya \texttt{RULDR}. \\
		\hline
		Total cost & Menyimpan akumulasi cost dari seluruh gerakan sliding yang sudah dilakukan. \\
		\hline
		Nilai heuristik & Digunakan oleh GBFS dan A* untuk memperkirakan kedekatan state terhadap target berikutnya. \\
		\hline
	\end{tabular}
\end{table}

Secara umum, identitas sebuah state ditentukan oleh kombinasi posisi aktor dan progres angka yang sudah dikumpulkan. 
Dua state dengan posisi aktor sama belum tentu setara apabila angka yang sudah dikumpulkan berbeda. 
Sebagai contoh, aktor yang berada pada posisi yang sama tetapi sudah mengumpulkan angka \texttt{0} memiliki kondisi pencarian yang berbeda dengan aktor yang belum mengumpulkan angka tersebut.

\section{Mekanisme Ekspansi State}

Pada setiap state, algoritma mencoba empat arah gerakan, yaitu atas, bawah, kiri, dan kanan. 
Arah tersebut direpresentasikan dengan karakter \texttt{U}, \texttt{D}, \texttt{L}, dan \texttt{R}. 
Setiap gerakan diproses menggunakan mekanisme sliding sampai aktor berhenti, keluar papan, melewati lava, atau melanggar aturan angka.

\begin{table}[H]
	\centering
	\caption{Aturan ekspansi gerakan}
	\begin{tabular}{|p{3.5cm}|p{10.5cm}|}
		\hline
		\rowcolor{gray!20}
		\textbf{Kondisi} & \textbf{Penanganan} \\
		\hline
		Menabrak \texttt{X} & Aktor berhenti tepat sebelum rintangan. Tile \texttt{X} tidak dilewati dan cost-nya tidak dihitung. \\
		\hline
		Melewati \texttt{L} & Gerakan dianggap menyebabkan \textit{game over}, sehingga state hasil tidak diproses. \\
		\hline
		Keluar papan & Jika aktor bergerak keluar batas papan tanpa tertahan rintangan, gerakan dianggap tidak valid. \\
		\hline
		Melewati angka benar & Angka diterima jika nilainya sama dengan angka terakhir yang dikumpulkan ditambah satu. \\
		\hline
		Melewati angka salah & Gerakan dianggap tidak valid karena melanggar constraint urutan angka. \\
		\hline
		Berhenti di \texttt{O} & State menjadi goal hanya jika seluruh angka sudah dikumpulkan sesuai urutan. \\
		\hline
	\end{tabular}
\end{table}

\begin{tcolorbox}[
	colback=gray!5,
	colframe=black!60,
	title=\textbf{Skema Umum Ekspansi State},
	arc=2mm,
	boxrule=0.8pt
]
\begin{center}
\begin{forest}
for tree={
	draw,
	rounded corners,
	align=center,
	parent anchor=south,
	child anchor=north,
	l sep=8mm,
	s sep=8mm,
	font=\small
}
[State saat ini
	[Coba \texttt{U}]
	[Coba \texttt{D}]
	[Coba \texttt{L}]
	[Coba \texttt{R}]
]
\end{forest}
\end{center}

Setiap arah tidak langsung menghasilkan perpindahan satu petak, melainkan disimulasikan sebagai gerakan sliding sampai kondisi berhenti atau kondisi tidak valid ditemukan.
\end{tcolorbox}

\section{Uniform Cost Search}

Uniform Cost Search atau UCS adalah algoritma pencarian yang memilih state berdasarkan total cost aktual paling kecil. 
Pada permasalahan ini, UCS digunakan untuk mencari solusi dengan cost minimum, karena setiap tile memiliki nilai cost yang berbeda-beda. 
Dengan demikian, jalur yang memiliki jumlah gerakan lebih sedikit belum tentu merupakan jalur terbaik apabila melewati tile dengan cost besar.

\begin{tcolorbox}[
	colback=green!5,
	colframe=green!50!black,
	title=\textbf{Prinsip UCS},
	arc=2mm,
	boxrule=0.8pt
]
UCS selalu memilih state dengan nilai $g(n)$ terkecil.
\[
	g(n) = \text{total cost dari start sampai state } n
\]
Dengan strategi ini, UCS memprioritaskan jalur yang benar-benar paling murah berdasarkan cost yang sudah dikeluarkan.
\end{tcolorbox}

\subsection{Langkah Kerja UCS}

\begin{longtable}{|p{1.2cm}|p{12.8cm}|}
	\caption{Langkah kerja Uniform Cost Search} \\
	\hline
	\rowcolor{green!15}
	\textbf{No.} & \textbf{Langkah} \\
	\hline
	\endfirsthead
	
	\hline
	\rowcolor{green!15}
	\textbf{No.} & \textbf{Langkah} \\
	\hline
	\endhead
	
	1 & Membuat initial state dari posisi awal aktor dengan cost awal $0$, angka terakhir $-1$, dan path kosong. \\
	\hline
	2 & Memasukkan initial state ke dalam priority queue. \\
	\hline
	3 & Mengambil state dengan nilai $g(n)$ terkecil dari priority queue. \\
	\hline
	4 & Memeriksa apakah state tersebut sudah memenuhi kondisi goal. \\
	\hline
	5 & Jika belum goal, algoritma mengekspansi state ke empat arah: \texttt{U}, \texttt{D}, \texttt{L}, dan \texttt{R}. \\
	\hline
	6 & Untuk setiap arah, dilakukan simulasi sliding sampai aktor berhenti atau gerakan menjadi tidak valid. \\
	\hline
	7 & Jika gerakan valid, cost baru dihitung dengan menambahkan cost sliding ke cost state saat ini. \\
	\hline
	8 & State baru dimasukkan ke priority queue apabila belum pernah dicapai sebelumnya dengan cost yang lebih kecil atau sama. \\
	\hline
	9 & Proses diulang sampai solusi ditemukan atau priority queue kosong. \\
	\hline
\end{longtable}

\subsection{Struktur Data UCS}

\begin{table}[H]
	\centering
	\caption{Struktur data yang digunakan pada UCS}
	\begin{tabular}{|p{4cm}|p{10cm}|}
		\hline
		\rowcolor{green!15}
		\textbf{Struktur Data} & \textbf{Fungsi} \\
		\hline
		Priority queue & Menyimpan state yang akan diproses, dengan prioritas berdasarkan nilai $g(n)$ terkecil. \\
		\hline
		Map best cost & Menyimpan cost terbaik untuk setiap kombinasi posisi dan collected mask. \\
		\hline
		State & Menyimpan posisi aktor, path, cost, collected mask, dan angka terakhir. \\
		\hline
		MoveResult & Menyimpan hasil simulasi sliding, seperti posisi baru, cost gerakan, tile yang dilewati, dan status validitas gerakan. \\
		\hline
	\end{tabular}
\end{table}

\subsection{Kelebihan dan Kekurangan UCS}

\begin{table}[H]
	\centering
	\caption{Kelebihan dan kekurangan UCS}
	\begin{tabular}{|p{3cm}|p{11cm}|}
		\hline
		\rowcolor{green!15}
		\textbf{Aspek} & \textbf{Penjelasan} \\
		\hline
		Kelebihan & UCS dapat menghasilkan solusi optimal terhadap total cost selama semua cost bernilai positif. Hal ini cocok untuk puzzle yang memiliki variasi cost pada setiap tile. \\
		\hline
		Kekurangan & UCS dapat meninjau banyak konfigurasi karena tidak memiliki informasi estimasi arah menuju target. Pada papan besar atau konfigurasi kompleks, jumlah state yang dieksplorasi dapat menjadi sangat besar. \\
		\hline
	\end{tabular}
\end{table}

\section{Greedy Best-First Search}

Greedy Best-First Search atau GBFS adalah algoritma pathfinding yang memilih state berdasarkan nilai heuristik terkecil. 
Berbeda dengan UCS yang memperhatikan total cost aktual, GBFS lebih fokus pada state yang terlihat paling dekat dengan target berikutnya. 
Target berikutnya dapat berupa angka yang harus dikumpulkan selanjutnya atau titik tujuan \texttt{O} apabila seluruh angka sudah terkumpul.

\begin{tcolorbox}[
	colback=orange!5,
	colframe=orange!70!black,
	title=\textbf{Prinsip GBFS},
	arc=2mm,
	boxrule=0.8pt
]
GBFS selalu memilih state dengan nilai $h(n)$ terkecil.
\[
	h(n) = \text{estimasi jarak dari state } n \text{ menuju target berikutnya}
\]
Algoritma ini bersifat rakus karena hanya mempertimbangkan estimasi terdekat, bukan total cost aktual dari start.
\end{tcolorbox}

\subsection{Langkah Kerja GBFS}

\begin{longtable}{|p{1.2cm}|p{12.8cm}|}
	\caption{Langkah kerja Greedy Best-First Search} \\
	\hline
	\rowcolor{orange!20}
	\textbf{No.} & \textbf{Langkah} \\
	\hline
	\endfirsthead
	
	\hline
	\rowcolor{orange!20}
	\textbf{No.} & \textbf{Langkah} \\
	\hline
	\endhead
	
	1 & Membuat initial state dari posisi awal aktor. \\
	\hline
	2 & Menghitung nilai heuristik initial state terhadap target berikutnya. \\
	\hline
	3 & Memasukkan initial state ke priority queue dengan prioritas berdasarkan $h(n)$. \\
	\hline
	4 & Mengambil state dengan nilai heuristik terkecil. \\
	\hline
	5 & Memeriksa apakah state sudah memenuhi kondisi goal. \\
	\hline
	6 & Jika belum goal, state diekspansi ke empat arah gerakan. \\
	\hline
	7 & Untuk setiap arah, dilakukan simulasi sliding dengan tetap memperhatikan rintangan, lava, batas papan, dan urutan angka. \\
	\hline
	8 & Jika gerakan valid, state baru dihitung nilai heuristiknya. \\
	\hline
	9 & State baru dimasukkan ke priority queue apabila layak untuk diproses. \\
	\hline
	10 & Proses berlanjut sampai solusi ditemukan atau tidak ada state tersisa. \\
	\hline
\end{longtable}

\subsection{Heuristik pada GBFS}

Heuristik digunakan untuk memperkirakan kedekatan aktor terhadap target yang sedang dikejar. 
Jika masih ada angka yang belum dikumpulkan, maka target heuristik adalah angka berikutnya. 
Jika seluruh angka sudah dikumpulkan, maka target heuristik adalah titik tujuan \texttt{O}.

\begin{table}[H]
	\centering
	\caption{Peran heuristik pada GBFS}
	\begin{tabular}{|p{4cm}|p{10cm}|}
		\hline
		\rowcolor{orange!20}
		\textbf{Aspek} & \textbf{Penjelasan} \\
		\hline
		Target heuristik & Angka berikutnya yang harus dikumpulkan, atau \texttt{O} jika semua angka sudah terkumpul. \\
		\hline
		Dasar prioritas & State dengan nilai heuristik lebih kecil akan diproses lebih dahulu. \\
		\hline
		Hubungan dengan cost & Cost aktual tetap dihitung untuk total solusi, tetapi tidak menjadi prioritas utama dalam pemilihan state. \\
		\hline
		Hubungan dengan sliding & Heuristik hanya mengarahkan pencarian, sedangkan validitas gerakan tetap ditentukan oleh hasil simulasi sliding. \\
		\hline
	\end{tabular}
\end{table}

\begin{tcolorbox}[
	colback=orange!5,
	colframe=orange!70!black,
	title=\textbf{Ilustrasi Keputusan GBFS},
	arc=2mm,
	boxrule=0.8pt
]
\begin{center}
\begin{tabular}{c c c}
\textbf{State A} & \textbf{State B} & \textbf{State C} \\
$h(n)=8$ & $h(n)=3$ & $h(n)=5$ \\
\end{tabular}
\end{center}

GBFS akan memilih \textbf{State B} terlebih dahulu karena memiliki nilai heuristik paling kecil, meskipun cost aktual dari start menuju State B belum tentu paling kecil.
\end{tcolorbox}

\subsection{Kelebihan dan Kekurangan GBFS}

\begin{table}[H]
	\centering
	\caption{Kelebihan dan kekurangan GBFS}
	\begin{tabular}{|p{3cm}|p{11cm}|}
		\hline
		\rowcolor{orange!20}
		\textbf{Aspek} & \textbf{Penjelasan} \\
		\hline
		Kelebihan & GBFS cenderung lebih cepat menemukan solusi karena pencariannya diarahkan oleh heuristik menuju target berikutnya. \\
		\hline
		Kekurangan & GBFS tidak selalu menghasilkan solusi optimal karena tidak memprioritaskan total cost aktual. State yang tampak dekat ke target bisa saja melalui jalur dengan cost besar. \\
		\hline
	\end{tabular}
\end{table}

\section{A*}

A* adalah algoritma pathfinding yang menggabungkan prinsip UCS dan GBFS. 
Algoritma ini tidak hanya memperhatikan cost aktual dari start menuju state saat ini, tetapi juga memperhitungkan estimasi cost menuju target berikutnya. 
Dengan demikian, A* berusaha menyeimbangkan pencarian yang murah secara cost dan pencarian yang terarah menuju goal.

\begin{tcolorbox}[
	colback=red!5,
	colframe=red!60!black,
	title=\textbf{Prinsip A*},
	arc=2mm,
	boxrule=0.8pt
]
A* memilih state berdasarkan nilai $f(n)$ terkecil.
\[
	f(n) = g(n) + h(n)
\]
dengan:
\[
	g(n) = \text{cost aktual dari start ke state } n
\]
\[
	h(n) = \text{estimasi cost dari state } n \text{ ke target berikutnya}
\]
\end{tcolorbox}

\subsection{Langkah Kerja A*}

\begin{longtable}{|p{1.2cm}|p{12.8cm}|}
	\caption{Langkah kerja A*} \\
	\hline
	\rowcolor{red!15}
	\textbf{No.} & \textbf{Langkah} \\
	\hline
	\endfirsthead
	
	\hline
	\rowcolor{red!15}
	\textbf{No.} & \textbf{Langkah} \\
	\hline
	\endhead
	
	1 & Membuat initial state dengan $g(n)=0$ dan path kosong. \\
	\hline
	2 & Menghitung nilai heuristik $h(n)$ untuk initial state. \\
	\hline
	3 & Menghitung nilai $f(n)=g(n)+h(n)$ dan memasukkan state ke priority queue. \\
	\hline
	4 & Mengambil state dengan nilai $f(n)$ terkecil dari priority queue. \\
	\hline
	5 & Memeriksa apakah state sudah memenuhi kondisi goal. \\
	\hline
	6 & Jika belum goal, state diekspansi ke empat arah gerakan. \\
	\hline
	7 & Untuk setiap gerakan valid, algoritma menghitung cost sliding dan memperbarui nilai $g(n)$. \\
	\hline
	8 & Algoritma menghitung nilai $h(n)$ untuk state baru berdasarkan target berikutnya. \\
	\hline
	9 & Nilai $f(n)$ state baru dihitung dari $g(n)+h(n)$. \\
	\hline
	10 & State baru diproses hanya jika cost-nya lebih baik daripada kunjungan sebelumnya pada posisi dan collected mask yang sama. \\
	\hline
	11 & Proses diulang sampai solusi ditemukan atau seluruh kemungkinan state habis. \\
	\hline
\end{longtable}

\subsection{Peran $g(n)$, $h(n)$, dan $f(n)$}

\begin{table}[H]
	\centering
	\caption{Definisi fungsi biaya pada A*}
	\begin{tabular}{|p{3cm}|p{11cm}|}
		\hline
		\rowcolor{red!15}
		\textbf{Fungsi} & \textbf{Makna} \\
		\hline
		$g(n)$ & Total cost aktual dari initial state menuju state $n$. Nilainya bertambah berdasarkan cost seluruh tile yang dilewati selama sliding. \\
		\hline
		$h(n)$ & Estimasi cost dari state $n$ menuju target berikutnya. Target dapat berupa angka berikutnya atau titik tujuan \texttt{O}. \\
		\hline
		$f(n)$ & Nilai prioritas state pada A*, yaitu gabungan antara cost aktual dan estimasi menuju target. \\
		\hline
	\end{tabular}
\end{table}

\subsection{Struktur Data A*}

\begin{table}[H]
	\centering
	\caption{Struktur data yang digunakan pada A*}
	\begin{tabular}{|p{4cm}|p{10cm}|}
		\hline
		\rowcolor{red!15}
		\textbf{Struktur Data} & \textbf{Fungsi} \\
		\hline
		Priority queue & Menyimpan state yang akan diproses dengan prioritas berdasarkan nilai $f(n)$ terkecil. \\
		\hline
		Map best cost & Menyimpan cost terbaik untuk setiap kombinasi posisi aktor dan collected mask. \\
		\hline
		State & Menyimpan posisi, path, $g(n)$, $h(n)$, collected mask, dan angka terakhir. \\
		\hline
		Heuristic & Menghitung estimasi jarak atau cost menuju target berikutnya. \\
		\hline
		Movement engine & Mensimulasikan gerakan sliding dan menentukan validitas hasil gerakan. \\
		\hline
	\end{tabular}
\end{table}

\subsection{Kelebihan dan Kekurangan A*}

\begin{table}[H]
	\centering
	\caption{Kelebihan dan kekurangan A*}
	\begin{tabular}{|p{3cm}|p{11cm}|}
		\hline
		\rowcolor{red!15}
		\textbf{Aspek} & \textbf{Penjelasan} \\
		\hline
		Kelebihan & A* lebih terarah dibanding UCS karena menggunakan heuristik, tetapi tetap mempertimbangkan cost aktual. Jika heuristik yang digunakan admissible, A* dapat menghasilkan solusi optimal. \\
		\hline
		Kekurangan & Performa A* sangat bergantung pada kualitas heuristik. Jika heuristik kurang informatif, A* dapat berperilaku mirip UCS dan tetap mengeksplorasi banyak state. \\
		\hline
	\end{tabular}
\end{table}

\section{Perbandingan UCS, GBFS, dan A*}

Ketiga algoritma memiliki strategi pemilihan state yang berbeda. 
UCS berfokus pada cost aktual, GBFS berfokus pada estimasi ke target, sedangkan A* menggabungkan keduanya. 
Perbandingan ini penting karena karakteristik puzzle dipengaruhi oleh dua hal utama, yaitu cost tile yang bervariasi dan posisi target yang harus dicapai secara berurutan.

\begin{table}[H]
	\centering
	\caption{Perbandingan algoritma pathfinding}
	\begin{tabular}{|p{3cm}|p{3.5cm}|p{3.5cm}|p{3.5cm}|}
		\hline
		\rowcolor{gray!25}
		\textbf{Aspek} & \textbf{UCS} & \textbf{GBFS} & \textbf{A*} \\
		\hline
		Dasar prioritas & $g(n)$ & $h(n)$ & $f(n)=g(n)+h(n)$ \\
		\hline
		Memperhatikan cost aktual & Ya & Tidak sebagai prioritas utama & Ya \\
		\hline
		Menggunakan heuristik & Tidak & Ya & Ya \\
		\hline
		Kecepatan pencarian & Cenderung lebih lambat & Cenderung lebih cepat & Bergantung kualitas heuristik \\
		\hline
		Optimalitas solusi & Optimal untuk cost positif & Tidak selalu optimal & Optimal jika heuristik admissible \\
		\hline
		Kelemahan utama & Eksplorasi state besar & Dapat memilih jalur mahal & Bergantung pada heuristik \\
		\hline
	\end{tabular}
\end{table}

\begin{tcolorbox}[
	colback=purple!5,
	colframe=purple!60!black,
	title=\textbf{Ringkasan Pemilihan Algoritma},
	arc=2mm,
	boxrule=0.8pt
]
\begin{itemize}
	\item UCS cocok digunakan ketika prioritas utama adalah mendapatkan solusi dengan total cost minimum.
	\item GBFS cocok digunakan ketika pencarian ingin diarahkan secepat mungkin menuju target, meskipun tidak menjamin solusi optimal.
	\item A* menjadi pendekatan yang lebih seimbang karena mempertimbangkan cost aktual dan estimasi menuju target.
\end{itemize}
\end{tcolorbox}